\newpage
\phantomsection
\label{prf:images-chains-monotone}
\LRAProofFor{prop:images-chains-monotone}

\begin{remark*}[Return]
\hyperref[prop:images-chains-monotone]{Return to Proposition}
\end{remark*}

\begin{proposition*}[Images of Chains Under Monotone Maps]
If $C$ is a chain and $f$ is order-preserving or order-reversing, then
$f(C)$ is a chain.
\end{proposition*}

\begin{proof}[Professional Standard Proof]
\LRAProofBodyStart
Take $u,v\in f(C)$. Then $u=f(c_1)$ and $v=f(c_2)$ for some
$c_1,c_2\in C$. Since $C$ is a chain, either $c_1\leq c_2$ or
$c_2\leq c_1$. An order-preserving map keeps that comparison between $u$ and
$v$; an order-reversing map reverses it. In either case $u$ and $v$ are
comparable.
\end{proof}

\begin{proof}[Detailed Learning Proof]
\LRAProofBodyStart
A chain is a set where every two elements can be compared. A monotone or
antitone map may collapse elements or reverse comparisons, but it does not
turn a comparable pair into an incomparable pair.
\end{proof}
\begin{remark*}[Proof structure]
\end{remark*}



\begin{dependencies}
\begin{itemize}
  \item \hyperref[prop:images-chains-monotone]{Proof target}
\end{itemize}
\end{dependencies}

\clearpage

